\chapter{Eisenbahnentwurfsregeln}
\section*{Motivation}
Trassierungsentwurf ist keine rein geometrische oder mathematische Aufgabe. Er
wird von Engineering-Regeln, Sicherheitsanforderungen und operativen
Bedingungen bestimmt, die aus nationalen Vorschriften, internationalen Normen
wie UIC, betreiberspezifischen Richtlinien wie denen der Deutschen Bahn und
empirischer Engineering-Praxis stammen. AIM integriert sie in Modellierung und
Optimierung, statt sie nur als externe Prüfungen zu behandeln.

\section{Regelkategorien}
\subsection{Geometrische Regeln}
Mindestradius, Alignment-Stetigkeit, Anforderungen an Übergangslängen und
Krümmungsmonotonie in Übergängen.

\subsection{Kinematische und dynamische Regeln}
Grenzen von Querbeschleunigung und Ruck, Verträglichkeit von Krümmung und
Überhöhung sowie geschwindigkeitsabhängige Anforderungen.

\subsection{Operative Regeln}
Zulässige Geschwindigkeiten, Sicherheitsreserven, Instandhaltungsaspekte und
Interoperabilitätsbedingungen.

\subsection{Topologische Regeln}
Bedingungen an Weichen und Kreuzungen, Parallelgleisabstand und
Netzkonnektivität.

\section{Geometrische Bedingungen}
\subsection{Mindestradius}
\begin{definition}[Mindestradiusbedingung]
\[
|\kappa(s)|\leq\kappa_{\max},\qquad\kappa_{\max}=\frac1{R_{\min}}.
\]
\end{definition}
\(R_{\min}\) hängt von Entwurfsgeschwindigkeit, Fahrzeugeigenschaften und
regulatorischen Normen ab.

\subsection{Stetigkeitsbedingungen}
Alignments müssen mindestens \(C^0\)-Stetigkeit der Position,
\(C^1\)-Stetigkeit der Richtung und typischerweise \(C^2\)-Stetigkeit der
Krümmung erfüllen. Für Hochgeschwindigkeitsstrecken kann höhere Glattheit
erforderlich sein.

\subsection{Übergangslänge}
\begin{definition}[Minimale Übergangslänge]
\[
L\geq L_{\min}(v,\kappa_1,\kappa_2).
\]
\end{definition}
Typische Formulierungen beziehen die zulässige Änderungsrate der
Querbeschleunigung und Grenzen der Überhöhungsentwicklung ein.

\section{Dynamische Bedingungen}
\subsection{Querbeschleunigung}
\begin{definition}[Querbeschleunigung]
\[a_y(s)=v^2\kappa(s).\]
\end{definition}
\begin{constraint}\[|a_y(s)|\leq a_{\max}.\]\end{constraint}

\subsection{Überhöhung und Überhöhungsfehlbetrag}
\begin{definition}[Überhöhungswinkel]
\(\phi(s)\) bezeichnet den Überhöhungswinkel.
\end{definition}
\begin{definition}[Effektive Beschleunigung]
\[
a_{\text{eff}}(s)=v^2\kappa(s)-g\tan(\phi(s)).
\]
\end{definition}
\begin{constraint}\[|a_{\text{eff}}(s)|\leq a_{\text{eff,max}}.\]\end{constraint}

\subsection{Ruck}
\begin{definition}[Ruck]
\[j(s)=\frac{\dd}{\dd t}a_y(s).\]
\end{definition}
Mit \(s\) als Parameter gilt \(j(s)=v^3\kappa'(s)\).
\begin{constraint}\[|j(s)|\leq j_{\max}.\]\end{constraint}

\section{Kopplung von Krümmung und Überhöhung}
\begin{constraint}
\[
|\kappa'(s)|\leq f_1(v,\phi),\qquad
|\phi'(s)|\leq f_2(v,\kappa).
\]
\end{constraint}
Krümmungsänderungen erfordern entsprechende Überhöhungsentwicklung;
Überhöhungsänderungen müssen dynamische Grenzen einhalten.

\section{Geschwindigkeitsabhängige Bedingungen}
\begin{definition}[Geschwindigkeitsprofil]
Ein Geschwindigkeitsprofil ist eine Funktion
\[
v:[0,L]\to\R_+.
\]
\end{definition}
Alle maßgeblichen Bedingungen müssen punktweise bezüglich \(v(s)\) bewertet
werden.

\section{Darstellung von Bedingungen in AIM}
\subsection{Abstrakte Form}
\begin{definition}[Bedingungsfunktional]
Eine Bedingung wird dargestellt als
\[
C_i[\kappa,\phi,v]\leq0.
\]
\end{definition}
Damit kann sie kontinuierlich bewertet, numerisch diskretisiert und in ein
einheitliches Bedingungssystem aufgenommen werden.

\subsection{Diskrete Darstellung}
Numerisch werden Bedingungen auf \(s_0,\dots,s_N\) bewertet.
\begin{definition}[Diskrete Bedingung]
\[
C_i(\mathbf x)=\max_k C_i(s_k).
\]
\end{definition}

\section{Integration in die Optimierung}
Bedingungen gehen als \(g_i(\mathbf x)\leq0\) ein und können als harte
Machbarkeitsbedingungen, Strafglieder oder Barrieremethoden behandelt werden.

\section{Bezug zu Normen}
Die dargestellten Bedingungen entsprechen Regeln aus UIC-Codes, nationalen
Eisenbahnvorschriften und Betreiberrichtlinien. AIM stellt eine mathematische
Darstellung bereit, die sie über verschiedene Datenquellen konsistent anwendet.

\section{Diskussion}
Traditionelle Abläufe wenden Regeln nach dem Entwurf als externe Prüfungen an.
AIM integriert sie unmittelbar in Modellierung, Rekonstruktion und Optimierung
und verschiebt damit Validierung zu bedingungsgetriebenem Entwurf.

\section{Zusammenfassung}
\begin{summary}
Eisenbahnentwurfsregeln definieren die zulässige Menge von
Alignment-Lösungen. AIM formalisiert sie als mathematische Bedingungen und
integriert sie in die Optimierung. Alignment-Entwurf wird damit zum
beschränkten Optimierungsproblem unter geometrischen und dynamischen
Anforderungen.
\end{summary}
